%%%%%%%%%%%%%%%%%%%%%%%
%%% DOCUMENT SETTINGS
%%%%%%%%%%%%%%%%%%%%%%%
\documentclass[11pt]{exam}
\printanswers % Uncomment to show answers
%%%%%%%%%%%%%%
%%% PACKAGES
%%%%%%%%%%%%%%
\usepackage{amsmath,amsfonts,amssymb,amsthm} % math fonts and symbols
\usepackage{bbm} % Math blackboard font
\usepackage{enumitem} % For reference enumerations lists
\usepackage{textcomp} % some symbols (like copyright)
\usepackage[top = 2cm, bottom = 2cm, right = 1.2cm, left = 1.2cm, paper = a4paper, headheight = 6cm]{geometry} % Margins setup
\usepackage{xcolor} % Colors
\usepackage{graphicx}
%%%%%%%%%%%%%%%%%%%%%%%%%%
%%% MACROS AND COMMANDS
%%%%%%%%%%%%%%%%%%%%%%%%%%
\newcommand{\N}{\mathbb{N}} % natural number set
\newcommand{\R}{\mathbb{R}} % real number set
\makeatletter
\renewcommand{\@makefnmark}{\textsuperscript{\arabic{footnote}}} % Ensure numeric superscripts
\renewcommand{\thefootnote}{\arabic{footnote}} % Ensure numeric labels
\makeatother
\setcounter{footnote}{1}
\newcommand{\BAR}{%
    \hspace{-\arraycolsep}%
    \strut\vrule % the `\vrule` is as high and deep as a strut
    \hspace{-\arraycolsep}%
}
%%%%%%%%%%%%%%%%%%%%%%%%%%%%%%%%%
%%% HEADER AND FOOT 
%%%%%%%%%%%%%%%%%%%%%%%%%%%%%%%%%
% Defining information
\newcommand{\degree}{Bachelor in Philosophy, Politics and Economics}
\newcommand{\shortdeg}{PPE}
\newcommand{\exam}{1st Midterm}
\newcommand{\subject}{Mathematics}
\newcommand{\theyear}{2024/25}
\newcommand{\ccauthor}{A. G. Azze}
\newcommand{\thedate}{Oct 17, 2024}
\newcommand{\university}{CUNEF Universidad}
% Setting header and footer
\pagestyle{headandfoot}
\header{\degree \\ \subject\ - \theyear}{}{\thedate \\ \ccauthor}
\footer{\university}{}{\thepage}
%%%%%%%%%%%%%%%%%%%%%%%%%
%%% DOCUMENT CONTENT
%%%%%%%%%%%%%%%%%%%%%%%%%
\begin{document}
%--- Instructions ---%
\begin{center}

    {\Huge \textbf{\exam{}}}
    \vspace{0.7cm}

    \begin{flushright}
        \textbf{Time Limit:} 50 minutes
    \end{flushright}
    \vspace{-0.2cm}
    
    \fbox{\parbox{\textwidth}{\centering 
        \begin{enumerate}[leftmargin = 0.5cm, label = $\bullet$, rightmargin = 0.2cm]
    
            \item \textbf{Every non-trivial calculation and claim in the exam must be properly justified.}

            \item The use of calculators is not necessary and therefore not allowed.

            \item Digital devices such as mobile phones, tablets, and laptops, as well as internet access, are forbidden.
            
        \end{enumerate}
    }}
\end{center}
%--- Instructions ---%

%--- Questions ---%

\begin{questions}

% Question
\question[3]
    Consider the function $f(x) = \frac{\sqrt{x + 11} - 4}{x - 5}$.
    \begin{enumerate}[leftmargin = 0.5cm, label = (\alph*), rightmargin = 0.2cm]
        \item What is the domain of $f$?
        \item What is the continuity domain of $f$?
        \item Calculate the limit $\lim_{x\rightarrow 5} f(x)$.
        \item For $h(x) = e^{x} + x^2$ and $g(x) = 8x$, obtain the limit $\lim_{x\rightarrow 5} (h\circ g \circ f)(x)$.
    \end{enumerate}
% Solution
\begin{solution}[0cm]
    \begin{enumerate}[leftmargin = 0.5cm, label = (\alph*), rightmargin = 0.2cm]
        
        \item Note that $f(x) = r_1(x)/r_2(x)$, with $r_1(x) = \sqrt{x + 11} - 4$ and $r_2(x) = x - 5$. Therefore, the domain of $f$, call it $D$, is the intersection of the domain of $g$ and the domain of $h$, subtracting the points where $r_2(h) = 0$. That is, $D = [-11, \infty)\backslash\{5\}$.
        
        \item Since $g(x)$ and $h(x)$ are continuous functions in their whole domain (they only involve adding constants and the square root function, which is continuous), then $f(x)$ is also a continuous function in its whole domain. That is, the continuity domain of $f$ is also $D$.
        
        \item The limit can be computed using the ``conjugate trick'':
        \begin{align*}
            \lim_{x\rightarrow 5} \frac{\sqrt{x + 11} - 4}{x - 5} 
            &= \lim_{x\rightarrow 5} \frac{\sqrt{x + 11} - 4}{x - 5}\frac{\sqrt{x + 11} + 4}{\sqrt{x + 11} + 4} 
            = \lim_{x\rightarrow 5} \frac{\sqrt{x + 11}^2 - 4^2}{(x - 5)\sqrt{x + 11} + 4} \\
            &= \lim_{x\rightarrow 5} \frac{1}{\sqrt{x + 11} + 4} 
            = 1/8.
        \end{align*}
        
        \item The functions $g(x)$ and $h(x)$ are continuous for all $x\in\R$, then
        \begin{align*}
            \lim_{x\rightarrow 5} h(g(f(x))) = h(g(\lim_{x\rightarrow 5} f(x))) = h(g(1/8)) = e + 1.
        \end{align*}
        
    \end{enumerate}
\end{solution}

% Question
\question[3]
    Consider the function $f(x) = e^{e^{x}}$.

    \begin{enumerate}[leftmargin = 0.5cm, label = (\alph*), rightmargin = 0.2cm]
        
        \item Compute the derivative of $f(x)$.

        \item Use the value of $f'(0)$ to offer an approximation of the function at $f(0.01)$.

        \item Are the limits $l_1 = \lim_{x\rightarrow 0} \frac{f(x) - e}{x}$ and $l_ 2 = \lim_{x \rightarrow\infty} xf(x)$ determinate or indeterminate forms of limits?

        \item Obtain the values of $l_1$ and $l_2$, if they exist.
        
    \end{enumerate}
    
%% Solution
\begin{solution}[0cm]
    
    \begin{enumerate}[leftmargin = 0.5cm, label = (\alph*), rightmargin = 0.2cm]
        
        \item Note that $f(x) = g(h(x))$, for $g(x) = e^x$ and $h(x) = x^2$. Then, by using the composition rule twice, we obtain that
            \begin{align*}
                f'(x) = h'(x)g'(h(x)) = e^{x}e^{e^{x}}.    
            \end{align*}
            
        \item We know that $f(a + \Delta) \approx f(a) + f'(a)\Delta$. Taking $a = 0$ and $\Delta = 0.01$, we obtain that $f(0.01) \approx f(0) + f'(0)0.01$. From the previous part we get that $f'(0) = e$, and from the definition of $f$ we get that $f(0) = e$. Hence,
        \begin{align*}
            f(0.01) \approx e(1 + 0.01) \approx 2.74546.
        \end{align*}
        The actual value of $f(0.01)$ is $e^{e^{0.01}} \approx 2.74574$, for you to appreciate the accuracy of the estimation.

        \item The limit $l_2$ is a determinate form of the type $\infty\times\infty$, while the limit $l_1$ is an indeterminate form of the type $0/0$. 

        \item It is trivial to get $l_2 = \infty$. Using L'Hospital rule we obtain that
        \begin{align*}
            l_1 = \lim_{x\rightarrow 0} \frac{f(x) - e}{x^2} 
            = \lim_{x\rightarrow 0} f'(x)
            = \lim_{x\rightarrow 0} e^{x}e^{e^{x}}
            = e.
        \end{align*}
        
    \end{enumerate}
    
    
\end{solution}

% Question
\question[1]
    Let $f(x)$ and $g(x)$ be an increasing and decreasing function, respectively. Is the function $h(x) = g(f(x))$ increasing, decreasing, or neither? Justify your answer with a proof.
%% Solution
\begin{solution}[0cm]
    Take $x_1 \leq x_2$; then, since $f$ is increasing, $f(x_1) \geq f(x_2)$; then since $g$ is decreasing, $h(x_1) = g(f(x_1)) \leq g(f(x_2))$, meaning that $h$ is decreasing.
\end{solution}

% Question
\question[3]
    You own a barrel of wine whose initial price in the market is $K = 100$ dollars. As the wine gets older, its price increases, such that the barrel price after $t$ years is given by $W(t) = e^{2\sqrt{t}}K$. On the other hand, the inflation rate of the country makes future profits decrease exponentially at a rate of $0.5$, that is, $x$ dollars today equals $e^{-0.5t}x$ dollars $t$ years from now. 

    \begin{enumerate}[leftmargin = 0.5cm, label = (\alph*), rightmargin = 0.2cm]
        
        \item Obtain the function $P(t)$ that determines the wine barrel price after $t$ years

        \item When should you sell the wine to maximize your profit and what would be that optimal profit?

        \item Prove that the optimal time to sell does not change if the value of $K$ changes.
        
    \end{enumerate}

% Solution
\begin{solution}[0cm]
    \begin{enumerate}[leftmargin = 0.5cm, label = (\alph*), rightmargin = 0.2cm]
        
        \item The wine barrel price after $t$ years is given by $P(t) = Ke^{2\sqrt{t} - 0.5t}$.

        \item Note that $P'(t) = Ke^{2\sqrt{t} - 0.5t}\left(\frac{1}{\sqrt{t}} - 0.5\right)$. Then, $P'(t) > 0$ for all $0 < t < 4$, $P'(4) = 0$, and $P'(t) < 0$ for all $t > 4$, meaning that the function increases until $t = 4$, where it reaches a maximum, and decreases afterwards. The maximum profit is given $P(4) = 100e^{2\sqrt{4} - 2} = 100e^2$.

        \item Note that the point of maximum ($t = 4$) obtained in the previous part does not depend on $K$, hence, The profit made by considering an arbitrary $K$ gets also maximized at $t = 4$.
        
    \end{enumerate}
\end{solution}

\end{questions}	

\begin{center}
You may find useful the following derivatives of basic functions and rules of derivatives
\end{center}

\begin{center}
    \textbf{Differentiation rules}
\end{center}
        
        \begin{enumerate}[leftmargin = 0.5cm, label = , rightmargin = 0.2cm]

        \item \textbf{Linearity:} $h(x) = a f(x) + b g(x) \Longrightarrow h'(x) = a f'(x) + b g'(x)$, for $a, b \in\R$

        \item \textbf{Product rule:} $h(x) = f(x)g(x) \Longrightarrow h'(x) = f'(x)g(x) + f(x)g'(x)$

        \item \textbf{Quotient rule:} $h(x) = f(x)/g(x) \text{ and } g'(x) \neq 0 \Longrightarrow h'(x) = \frac{f'(x)g(x) - f(x)g'(x)}{g(x)^2}$

        \item \textbf{Chain rule:} $h(x) = g(f(x)) \Longrightarrow h'(x) = g'(f(x))f'(x)$

        \item \textbf{Inverse rule:} $h(x) = f^{-1}(y) \Longrightarrow h'(y) = 1/f'(f^{-1}(y))$
        
        \end{enumerate}

\begin{center}
    \textbf{Derivative of usual functions}
\end{center}

        \begin{enumerate}[leftmargin = 0.5cm, label = , rightmargin = 0.2cm]

        \item \textbf{Constant:} $f(x) = a,\; a \in\R \Longrightarrow f'(x) = 0$
        
        \item \textbf{Power:} $f(x) = x^a \Longrightarrow f'(x) = a x^{a-1}$
        

        \item \textbf{Exponential:} $f(x) = e^x \Longrightarrow f'(x) = e^x$

        \item \textbf{Logarithm:} $f(x) = \ln(x) \Longrightarrow f'(x) = 1/x$

        \item \textbf{Trigonometric:} $f(x) = \sin(x),\; g(x) = \cos(x) \Longrightarrow f'(x) = \cos(x),\; g'(x) = -\sin(x)$
        
        \end{enumerate}

\end{document}